\documentclass[12pt,a4paper]{article}

% --- Metadata (per course) ---
\def\courseshortheader{THỐNG KÊ ĐA BIẾN}
\def\coverbookline{NHÓM 5 --- XÁC SUẤT VÀ THỐNG KÊ}
\def\covermaintitle{Thống kê đa biến}
\def\coversubtitle{Giáo án lý thuyết --- PCA, nhân tố, phân biệt, phân cụm}

% Shared preamble for nhom_5 (requires \courseshortheader before \input)
\usepackage[utf8]{inputenc}
\usepackage[vietnamese]{babel}
\usepackage{amsmath,amssymb,amsthm}
\usepackage{geometry}
\usepackage{enumitem}
\usepackage[most]{tcolorbox}
\usepackage{hyperref}
\usepackage{fancyhdr}
\usepackage{graphicx}
\usepackage{booktabs}
\usepackage{tikz}
\usetikzlibrary{calc,arrows.meta,decorations.pathreplacing,positioning}
\usepackage{eso-pic}
\usepackage{xcolor}

\geometry{a4paper, margin=2cm, bottom=2.5cm}
\hypersetup{colorlinks=true, linkcolor=blue!70!black, urlcolor=blue}
\setlist[itemize]{topsep=4pt, itemsep=2pt}
\setlist[enumerate]{topsep=4pt, itemsep=2pt}

\AddToShipoutPictureFG{%
  \AtPageCenter{%
    \begin{tikzpicture}[remember picture, overlay]
      \node[rotate=45, scale=1, opacity=0.06]{%
        \fontsize{60}{72}\selectfont\bfseries\sffamily Đặng Tấn Quí%
      };
    \end{tikzpicture}%
  }%
}

\pagestyle{fancy}
\fancyhf{}
\fancyhead[L]{\small\textit{\courseshortheader}}
\fancyhead[R]{\small\textit{Đặng Tấn Quí}}
\fancyfoot[C]{\thepage}
\fancyfoot[L]{\footnotesize\textcopyright\ Đặng Tấn Quí}
\fancyfoot[R]{\footnotesize SĐT: 0377\,122\,210}
\renewcommand{\headrulewidth}{0.4pt}
\renewcommand{\footrulewidth}{0.4pt}

\fancypagestyle{plain}{%
  \fancyhf{}
  \fancyfoot[C]{\thepage}
  \fancyfoot[L]{\footnotesize\textcopyright\ Đặng Tấn Quí}
  \fancyfoot[R]{\footnotesize SĐT: 0377\,122\,210}
  \renewcommand{\headrulewidth}{0pt}
  \renewcommand{\footrulewidth}{0.4pt}
}

\newtcolorbox{dongco}[1][]{%
  enhanced, breakable,
  colback=teal!4!white, colframe=teal!60!black,
  fonttitle=\bfseries, title={Động cơ học -- Tại sao cần khái niệm này?}, #1%
}
\newtcolorbox{ynghia}[1][]{%
  enhanced, breakable,
  colback=purple!3!white, colframe=purple!50!black,
  fonttitle=\bfseries, title={Ý nghĩa \& Trực giác}, #1%
}
\newtcolorbox{ungdung}[1][]{%
  enhanced, breakable,
  colback=olive!5!white, colframe=olive!60!black,
  fonttitle=\bfseries, title={Ứng dụng thực tế}, #1%
}
\newtcolorbox{dinhnghia}[1][]{%
  enhanced, breakable,
  colback=blue!3!white, colframe=blue!60!black,
  fonttitle=\bfseries, title={Định nghĩa}, #1%
}
\newtcolorbox{dinhly}[1][]{%
  enhanced, breakable,
  colback=green!3!white, colframe=green!50!black,
  fonttitle=\bfseries, title={Định lý}, #1%
}
\newtcolorbox{tinhchat}[1][]{%
  enhanced, breakable,
  colback=yellow!5!white, colframe=orange!50!black,
  fonttitle=\bfseries, title={Tính chất}, #1%
}
\newtcolorbox{phuongphap}[1][]{%
  enhanced, breakable,
  colback=gray!5!white, colframe=gray!50!black,
  fonttitle=\bfseries, title={Phương pháp}, #1%
}
\newtcolorbox{luuy}[1][]{%
  enhanced, breakable,
  colback=red!3!white, colframe=red!50!black,
  fonttitle=\bfseries, title={Lưu ý}, #1%
}
\newtcolorbox{congthuc}[1][]{%
  enhanced, breakable,
  colback=cyan!3!white, colframe=cyan!50!black,
  fonttitle=\bfseries, title={Công thức}, #1%
}
\newtcolorbox{vidu}[1][]{%
  enhanced, breakable,
  colback=violet!3!white, colframe=violet!50!black,
  fonttitle=\bfseries, title={Ví dụ}, #1%
}
\newtcolorbox{phanvidu}[1][]{%
  enhanced, breakable,
  colback=red!2!white, colframe=red!40!black,
  fonttitle=\bfseries, title={Phản ví dụ}, #1%
}
\newtcolorbox{tongket}[1][]{%
  enhanced, breakable,
  colback=blue!4!white, colframe=blue!70!black,
  fonttitle=\bfseries, title={Tổng kết chương}, #1%
}


\begin{document}

\thispagestyle{empty}
\begin{tikzpicture}[remember picture, overlay]
  \fill[blue!80!black] (current page.south west) rectangle (current page.north east);
  \fill[blue!60!cyan, opacity=0.8]
    (current page.south west) -- (current page.north west) --
    ([xshift=-3cm]current page.north east) -- (current page.south east) -- cycle;
  \foreach \r/\o in {6/0.05, 4.5/0.08, 3/0.12} {
    \fill[white, opacity=\o] ([xshift=2cm, yshift=-2cm]current page.north east) circle (\r cm);
  }
  \draw[white, line width=2pt] ([yshift=-5cm]current page.north west) ++(2cm,0) -- ++(17cm,0);
  \draw[white, line width=2pt] ([yshift=-16cm]current page.north west) ++(2cm,0) -- ++(17cm,0);
  \node[anchor=west, white, font=\fontsize{28}{34}\bfseries\selectfont]
    at ([xshift=3cm, yshift=-7.5cm]current page.north west) {\coverbookline};
  \node[anchor=west, white, font=\fontsize{20}{26}\selectfont]
    at ([xshift=3cm, yshift=-9cm]current page.north west) {\covermaintitle};
  \node[anchor=west, white, font=\fontsize{16}{22}\selectfont]
    at ([xshift=3cm, yshift=-10.2cm]current page.north west) {\coversubtitle};
  \node[anchor=west, white, font=\Large\bfseries]
    at ([xshift=3cm, yshift=-13cm]current page.north west) {Biên soạn:};
  \node[anchor=west, yellow!80!white, font=\fontsize{24}{30}\bfseries\selectfont]
    at ([xshift=3cm, yshift=-14.5cm]current page.north west) {ĐẶNG TẤN QUÍ};
  \node[anchor=west, white!90, font=\large]
    at ([xshift=3cm, yshift=-18cm]current page.north west) {SĐT: 0377\,122\,210};
  \node[anchor=south, white!80, font=\small]
    at ([yshift=1.5cm]current page.south) {Tài liệu có bản quyền --- Cấm sao chép khi chưa được sự đồng ý của tác giả};
  \node[anchor=south east, white, font=\Large\bfseries]
    at ([xshift=-2cm, yshift=3cm]current page.south east) {\the\year};
\end{tikzpicture}

\newpage
\tableofcontents
\newpage

\section*{Lời nói đầu}
\addcontentsline{toc}{section}{Lời nói đầu}

Tài liệu thuộc \textbf{Nhóm 5 --- Xác suất và thống kê}, biên soạn theo cùng triết lý với giáo án Đại số tuyến tính: học để \emph{hiểu} và \emph{dùng}, không chỉ để ghi nhớ công thức.

\begin{enumerate}
  \item \textbf{Động cơ trước định nghĩa.} Mỗi phần lớn mở bằng ô \textsf{\color{teal!60!black} Động cơ học} --- câu hỏi thực tế hoặc bài toán mẫu cần khái niệm đó.
  \item \textbf{Trực giác trước công thức.} Ô \textsf{\color{purple!50!black} Ý nghĩa \& Trực giác} giải thích ``công thức đang nói gì'' (xác suất = tần suất dài hạn, kỳ vọng = trung bình có trọng số, v.v.).
  \item \textbf{Ví dụ có kiểm tra.} Ô \textsf{\color{violet!50!black} Ví dụ} nên tự làm trước khi đọc lời giải; phần thống kê ứng dụng cần gắn dữ liệu số cụ thể.
  \item \textbf{Tổng kết cuối mỗi chương.} Ô \textsf{\color{blue!70!black} Tổng kết chương} liệt kê kết quả phải nhớ và liên kết sang phần sau.
  \item \textbf{Phân tầng theo môn.} X1--X3 là nền; X4--X5 nâng cao lý thuyết; X6--X8 chuyên sâu (đa biến, Bayes, quá trình ngẫu nhiên).
\end{enumerate}

\noindent\textbf{Cách đọc:} Động cơ $\to$ Định nghĩa / Định lý $\to$ Ví dụ $\to$ Ý nghĩa $\to$ Tổng kết. Với môn thống kê, luôn hỏi: \emph{mẫu là gì, tham số nào chưa biết, giả thuyết $H_0$ là gì}.

\newpage


\section{Mở đầu}

\begin{dongco}[title={Động cơ --- Mở đầu}]
Dữ liệu thường có nhiều biến cùng lúc. Thống kê đa biến tìm cấu trúc ẩn (thành phần chính, nhân tố), phân nhóm (cụm), hoặc phân loại (phân biệt).
\end{dongco}

\newpage

\section{Phân phối chuẩn nhiều chiều}

\begin{dongco}[title={Động cơ --- Phân phối chuẩn nhiều chiều}]
Nền xác suất cho vectơ ngẫu nhiên và ma trận hiệp phương sai.
\end{dongco}

\begin{dinhnghia}
  $\mathbf{X} = (X_1, \ldots, X_p)^T \sim \mathcal{N}_p(\boldsymbol{\mu}, \boldsymbol{\Sigma})$:
  \[
    f(\mathbf{x}) = \frac{1}{(2\pi)^{p/2}|\boldsymbol{\Sigma}|^{1/2}}\exp\!\left(-\frac{1}{2}(\mathbf{x}-\boldsymbol{\mu})^T\boldsymbol{\Sigma}^{-1}(\mathbf{x}-\boldsymbol{\mu})\right).
  \]
\end{dinhnghia}

\begin{tinhchat}
  \begin{itemize}
    \item Mọi tổ hợp tuyến tính $\mathbf{a}^T\mathbf{X} \sim \mathcal{N}(\mathbf{a}^T\boldsymbol{\mu}, \mathbf{a}^T\boldsymbol{\Sigma}\mathbf{a})$.
    \item Phân phối biên: $X_i \sim \mathcal{N}(\mu_i, \sigma_{ii})$.
    \item Không tương quan $\Leftrightarrow$ độc lập (đặc trưng của chuẩn).
    \item $(\mathbf{X}-\boldsymbol{\mu})^T\boldsymbol{\Sigma}^{-1}(\mathbf{X}-\boldsymbol{\mu}) \sim \chi^2(p)$.
  \end{itemize}
\end{tinhchat}

\begin{dinhly}[title={Phân phối Wishart}]
  Nếu $\mathbf{X}_i \sim \mathcal{N}_p(\boldsymbol{\mu}, \boldsymbol{\Sigma})$ iid thì:
  $\mathbf{W} = \sum_{i=1}^n (\mathbf{X}_i - \bar{\mathbf{X}})(\mathbf{X}_i - \bar{\mathbf{X}})^T \sim W_p(n-1, \boldsymbol{\Sigma})$.
\end{dinhly}

\begin{dinhly}[title={Hotelling $T^2$}]
  $T^2 = n(\bar{\mathbf{X}} - \boldsymbol{\mu}_0)^T \mathbf{S}^{-1}(\bar{\mathbf{X}} - \boldsymbol{\mu}_0)$. Dưới $H_0: \boldsymbol{\mu} = \boldsymbol{\mu}_0$:
  \[
    \frac{n - p}{(n-1)p}\,T^2 \sim F(p, n-p).
  \]
\end{dinhly}

%% ================================================================
%%  CHƯƠNG 2: PHÂN TÍCH THÀNH PHẦN CHÍNH (PCA)
%% ================================================================

\begin{tongket}[title={Tổng kết: Phân phối chuẩn nhiều chiều}]
\begin{enumerate}
  \item Nắm lại các \textbf{định nghĩa} và \textbf{công thức} cốt lõi trong mục ``Phân phối chuẩn nhiều chiều''.
  \item Tự làm lại ít nhất một ví dụ (nếu có) hoặc áp dụng công thức vào một tình huống số.
  \item Ghi chú điều kiện áp dụng (giả thiết độc lập, phân phối chuẩn, $n$ đủ lớn, v.v.).
\end{enumerate}
\end{tongket}

\section{Phân tích thành phần chính (PCA)}

\begin{dongco}[title={Động cơ --- Phân tích thành phần chính (PCA)}]
Giảm chiều nhưng giữ tối đa phương sai --- trực quan hóa và tiền xử lý.
\end{dongco}

\begin{dinhnghia}
  PCA tìm các tổ hợp tuyến tính $Y_k = \mathbf{a}_k^T \mathbf{X}$ sao cho $\text{Var}(Y_k)$ lớn nhất, các $Y_k$ không tương quan.
\end{dinhnghia}

\begin{dinhly}
  Thành phần chính thứ $k$ ứng với \textbf{trị riêng} $\lambda_k$ và \textbf{vectơ riêng} $\mathbf{a}_k$ của $\boldsymbol{\Sigma}$ (hoặc $\mathbf{S}$), $\lambda_1 \ge \lambda_2 \ge \cdots \ge \lambda_p \ge 0$:
  \[
    \text{Var}(Y_k) = \lambda_k, \qquad \text{Tỉ lệ giải thích} = \frac{\lambda_k}{\sum_{j=1}^p \lambda_j}.
  \]
\end{dinhly}

\begin{phuongphap}[title={Thực hành PCA}]
  \begin{enumerate}
    \item Chuẩn hóa dữ liệu (trung bình 0, phương sai 1) nếu các biến khác đơn vị.
    \item Chéo hóa ma trận hiệp phương sai (hoặc tương quan) mẫu.
    \item Chọn số thành phần: scree plot, tỉ lệ phương sai tích lũy $\ge 80\%$.
    \item Giải thích: xem loading (hệ số trong $\mathbf{a}_k$) để hiểu ý nghĩa thành phần.
  \end{enumerate}
\end{phuongphap}

%% ================================================================
%%  CHƯƠNG 3: PHÂN TÍCH NHÂN TỐ (FA)
%% ================================================================

\begin{tongket}[title={Tổng kết: Phân tích thành phần chính (PCA)}]
\begin{enumerate}
  \item Nắm lại các \textbf{định nghĩa} và \textbf{công thức} cốt lõi trong mục ``Phân tích thành phần chính (PCA)''.
  \item Tự làm lại ít nhất một ví dụ (nếu có) hoặc áp dụng công thức vào một tình huống số.
  \item Ghi chú điều kiện áp dụng (giả thiết độc lập, phân phối chuẩn, $n$ đủ lớn, v.v.).
\end{enumerate}
\end{tongket}

\section{Phân tích nhân tố (Factor Analysis)}

\begin{dongco}[title={Động cơ --- Phân tích nhân tố (Factor Analysis)}]
Mô hình hóa biến quan sát bằng ít nhân tố tiềm ẩn.
\end{dongco}

\begin{dinhnghia}
  Mô hình: $\mathbf{X} - \boldsymbol{\mu} = \mathbf{L}\mathbf{F} + \boldsymbol{\varepsilon}$, với $\mathbf{F}$ ($m$ nhân tố chung), $\boldsymbol{\varepsilon}$ (nhiễu riêng), $\mathbf{L}$ (ma trận tải $p \times m$).

  $\boldsymbol{\Sigma} = \mathbf{L}\mathbf{L}^T + \boldsymbol{\Psi}$, $\boldsymbol{\Psi} = \text{diag}(\psi_1, \ldots, \psi_p)$.
\end{dinhnghia}

\begin{phuongphap}[title={Ước lượng và xoay}]
  \begin{itemize}
    \item Ước lượng $\mathbf{L}$: phương pháp thành phần chính, hoặc MLE.
    \item \textbf{Xoay nhân tố}: Varimax (trực giao), Promax (xiên) --- dễ giải thích hơn.
    \item \textbf{Điểm nhân tố}: ước lượng $\hat{\mathbf{F}}$ cho từng quan sát.
  \end{itemize}
\end{phuongphap}

%% ================================================================
%%  CHƯƠNG 4: PHÂN TÍCH PHÂN BIỆT VÀ PHÂN CỤM
%% ================================================================

\begin{tongket}[title={Tổng kết: Phân tích nhân tố (Factor Analysis)}]
\begin{enumerate}
  \item Nắm lại các \textbf{định nghĩa} và \textbf{công thức} cốt lõi trong mục ``Phân tích nhân tố (Factor Analysis)''.
  \item Tự làm lại ít nhất một ví dụ (nếu có) hoặc áp dụng công thức vào một tình huống số.
  \item Ghi chú điều kiện áp dụng (giả thiết độc lập, phân phối chuẩn, $n$ đủ lớn, v.v.).
\end{enumerate}
\end{tongket}

\section{Phân tích phân biệt và phân cụm}

\begin{dongco}[title={Động cơ --- Phân tích phân biệt và phân cụm}]
Hai hướng: có nhãn (phân loại) vs không nhãn (gom cụm).
\end{dongco}

\subsection{Phân tích phân biệt (Discriminant Analysis)}

\begin{dinhnghia}[title={LDA (Linear Discriminant Analysis)}]
  Phân loại $\mathbf{x}$ vào nhóm $k$ cực đại hóa:
  \[
    \delta_k(\mathbf{x}) = \mathbf{x}^T\boldsymbol{\Sigma}^{-1}\boldsymbol{\mu}_k - \frac{1}{2}\boldsymbol{\mu}_k^T\boldsymbol{\Sigma}^{-1}\boldsymbol{\mu}_k + \ln\pi_k,
  \]
  với $\pi_k = P(\text{nhóm } k)$. Giả sử $\boldsymbol{\Sigma}$ chung.
\end{dinhnghia}

\begin{dinhnghia}[title={QDA (Quadratic DA)}]
  Nới lỏng: mỗi nhóm có $\boldsymbol{\Sigma}_k$ riêng. Biên phân loại là bậc hai.
\end{dinhnghia}

\subsection{Phân cụm (Clustering)}

\begin{phuongphap}[title={K-means}]
  \begin{enumerate}
    \item Chọn $K$ tâm ban đầu.
    \item Gán mỗi điểm vào cụm gần nhất.
    \item Cập nhật tâm = trung bình cụm.
    \item Lặp đến hội tụ. Cực tiểu $\sum_{k=1}^K \sum_{i \in C_k} \|\mathbf{x}_i - \boldsymbol{\mu}_k\|^2$.
  \end{enumerate}
\end{phuongphap}

\begin{phuongphap}[title={Phân cụm phân cấp}]
  \begin{itemize}
    \item \textbf{Agglomerative} (bottom-up): ghép dần.
    \item \textbf{Divisive} (top-down): tách dần.
    \item Liên kết: single, complete, average, Ward.
    \item Kết quả: dendrogram.
  \end{itemize}
\end{phuongphap}

\begin{phuongphap}[title={Chọn số cụm}]
  \begin{itemize}
    \item Phương pháp khuỷu tay (elbow method).
    \item Silhouette score.
    \item Gap statistic.
  \end{itemize}
\end{phuongphap}

\begin{tongket}[title={Tổng kết: Phân tích phân biệt và phân cụm}]
\begin{enumerate}
  \item Nắm lại các \textbf{định nghĩa} và \textbf{công thức} cốt lõi trong mục ``Phân tích phân biệt và phân cụm''.
  \item Tự làm lại ít nhất một ví dụ (nếu có) hoặc áp dụng công thức vào một tình huống số.
  \item Ghi chú điều kiện áp dụng (giả thiết độc lập, phân phối chuẩn, $n$ đủ lớn, v.v.).
\end{enumerate}
\end{tongket}



\end{document}
